%%%%%%%%%%%%%%%%%
% This is an sample CV template created using altacv.cls
% (v1.7.2, 28 August 2024) written by LianTze Lim (liantze@gmail.com). Compiles with pdfLaTeX, XeLaTeX and LuaLaTeX.
%
%% It may be distributed and/or modified under the
%% conditions of the LaTeX Project Public License, either version 1.3
%% of this license or (at your option) any later version.
%% The latest version of this license is in
%%    http://www.latex-project.org/lppl.txt
%% and version 1.3 or later is part of all distributions of LaTeX
%% version 2003/12/01 or later.
%%%%%%%%%%%%%%%%

%% Use the "normalphoto" option if you want a normal photo instead of cropped to a circle
% \documentclass[10pt,a4paper,normalphoto]{altacv}

\documentclass[10pt,a4paper,ragged2e,withhyper]{altacv}
%% AltaCV uses the fontawesome5 and simpleicons packages.
%% See http://texdoc.net/pkg/fontawesome5 and http://texdoc.net/pkg/simpleicons for full list of symbols.

% Change the page layout if you need to
\geometry{left=1.25cm,right=1.25cm,top=1.25cm,bottom=1cm,columnsep=1.2cm}

% The paracol package lets you typeset columns of text in parallel
\usepackage{paracol}
\usepackage[english, russian]{babel}

% Change the font if you want to, depending on whether
% you're using pdflatex or xelatex/lualatex
% WHEN COMPILING WITH XELATEX PLEASE USE
% xelatex -shell-escape -output-driver="xdvipdfmx -z 0" sample.tex
\iftutex 
  % If using xelatex or lualatex:
  \setmainfont{Roboto Slab}
  \setsansfont{Lato}
  \renewcommand{\familydefault}{\sfdefault}
\else
  % If using pdflatex:
  \usepackage[rm]{roboto}
  \usepackage[defaultsans]{lato}
  % \usepackage{sourcesanspro}
  \renewcommand{\familydefault}{\sfdefault}
\fi

% Change the colours if you want to
\definecolor{SlateGrey}{HTML}{2E2E2E}
\definecolor{LightGrey}{HTML}{666666}
\definecolor{DarkPastelRed}{HTML}{450808}
\definecolor{PastelRed}{HTML}{8F0D0D}
\definecolor{GoldenEarth}{HTML}{E7D192}
\colorlet{name}{black}
\colorlet{tagline}{PastelRed}
\colorlet{heading}{DarkPastelRed}
\colorlet{headingrule}{GoldenEarth}
\colorlet{subheading}{PastelRed}
\colorlet{accent}{PastelRed}
\colorlet{emphasis}{SlateGrey}
\colorlet{body}{LightGrey}

% Change some fonts, if necessary
\renewcommand{\namefont}{\Huge\rmfamily\bfseries\iftutex\addfontfeatures{Kerning=Off}\fi}
\renewcommand{\personalinfofont}{\footnotesize}
\renewcommand{\cvsectionfont}{\Large\rmfamily\bfseries\iftutex\addfontfeatures{Kerning=Off}\fi}
\renewcommand{\cvsubsectionfont}{\large\bfseries}


% Change the bullets for itemize and rating marker
% for \cvskill if you want to
\renewcommand{\cvItemMarker}{{\small\textbullet}}
\renewcommand{\cvRatingMarker}{\faCircle}
% ...and the markers for the date/location for \cvevent
% \renewcommand{\cvDateMarker}{\faCalendar*[regular]}
% \renewcommand{\cvLocationMarker}{\faMapMarker*}


% If your CV/résumé is in a language other than English,
% then you probably want to change these so that when you
% copy-paste from the PDF or run pdftotext, the location
% and date marker icons for \cvevent will paste as correct
% translations. For example Spanish:
% \renewcommand{\locationname}{Ubicación}
% \renewcommand{\datename}{Fecha}


%% Use (and optionally edit if necessary) this .tex if you
%% want to use an author-year reference style like APA(6)
%% for your publication list
% % When using APA6 if you need more author names to be listed
% because you're e.g. the 12th author, add apamaxprtauth=12
\usepackage[backend=biber,style=apa6,sorting=ydnt]{biblatex}
\defbibheading{pubtype}{\cvsubsection{#1}}
\renewcommand{\bibsetup}{\vspace*{-\baselineskip}}
\AtEveryBibitem{%
  \makebox[\bibhang][l]{\itemmarker}%
  \iffieldundef{doi}{}{\clearfield{url}}%
}
\setlength{\bibitemsep}{0.25\baselineskip}
\setlength{\bibhang}{1.25em}


%% Use (and optionally edit if necessary) this .tex if you
%% want an originally numerical reference style like IEEE
%% for your publication list
\usepackage[backend=biber,style=ieee,sorting=ydnt,defernumbers=true]{biblatex}
%% For removing numbering entirely when using a numeric style
\setlength{\bibhang}{1.25em}
\DeclareFieldFormat{labelnumberwidth}{\makebox[\bibhang][l]{\itemmarker}}
\setlength{\biblabelsep}{0pt}
\defbibheading{pubtype}{\cvsubsection{#1}}
\renewcommand{\bibsetup}{\vspace*{-\baselineskip}}
\AtEveryBibitem{%
  \iffieldundef{doi}{}{\clearfield{url}}%
}


% Tighter section spacing
\renewcommand{\cvsection}[1]{%
  \nointerlineskip\medskip%
  {\color{heading}\cvsectionfont\MakeUppercase{#1}}\\[-1ex]%
  {\color{headingrule}\rule{\linewidth}{2pt}\par}\smallskip
}

% Tighter lists
\setlist{itemsep=0.5\baselineskip,after=\vspace{0.1\baselineskip}}

%% sample.bib contains your publications
\addbibresource{sample.bib}
% \usepackage{academicons}\let\faOrcid\aiOrcid
\begin{document}
\name{Рауль Алимбеков}
\tagline{Embedded разработчик}
%% You can add multiple photos on the left or right
\photoR{2.8cm}{photo}
% \photoL{2.5cm}{Yacht_High,Suitcase_High}

\personalinfo{%
  % Not all of these are required!
  \email{raulalimbekov@gmail.com}
  \phone{+7-987-626-7271}
  \location{Москва, Россия}
  \github{RaulIQ}
  \printinfo{\faTelegram}{@Raul\_no\_Rules}[https://t.me/Raul_no_Rules]

  %% You can add your own arbitrary detail with
  %% \printinfo{symbol}{detail}[optional hyperlink prefix]
  % \printinfo{\faPaw}{Hey ho!}[https://example.com/]

  %% Or you can declare your own field with
  %% \NewInfoFiled{fieldname}{symbol}[optional hyperlink prefix] and use it:
  % \NewInfoField{gitlab}{\faGitlab}[https://gitlab.com/]
  % \gitlab{your_id}
  %%
  %% For services and platforms like Mastodon where there isn't a
  %% straightforward relation between the user ID/nickname and the hyperlink,
  %% you can use \printinfo directly e.g.
  % \printinfo{\faMastodon}{@username@instace}[https://instance.url/@username]
  %% But if you absolutely want to create new dedicated info fields for
  %% such platforms, then use \NewInfoField* with a star:
  % \NewInfoField*{mastodon}{\faMastodon}
  %% then you can use \mastodon, with TWO arguments where the 2nd argument is
  %% the full hyperlink.
  % \mastodon{@username@instance}{https://instance.url/@username}
}

\makecvheader
%% Depending on your tastes, you may want to make fonts of itemize environments slightly smaller
\AtBeginEnvironment{itemize}{\small}

%% Set the left/right column width ratio.
\columnratio{0.65}

% Start a 2-column paracol. Both the left and right columns will automatically
% break across pages if things get too long.
\begin{paracol}{2}

  \cvsection{Опыт работы}

  \cvevent{Программист встраиваемых систем}{ЮВС Авиа (концерн «Калашников») · Сентябрь 2025 — настоящее время}{}{}

  \begin{itemize}
    \item \textbf{Система управления батареей (BMS)} (\textbf{C++}, ESP32-S3, \textbf{FreeRTOS}, собственные платы): разработал логику заряда (BQ25792), мониторинга состояния (BQ40Z50) и питания нагрузки (TPS55288), драйверы по datasheet; реализовал \textbf{SMBus} поверх \textbf{I2C multi-master} для доступа полётного контроллера к данным батареи и режим энергосбережения. Участвовал в \textbf{bring-up} плат, настраивал параметры BQ40Z50 в bqStudio.
    \item Разработал \textbf{ППРЧ} для связи дрона с наземной станцией по \textbf{UDP} на \textbf{Embedded Linux} (Raspberry Pi, wfb-ng, RTL8812EU): 10 прыжков/с по 42 каналам, псевдослучайная последовательность на SHA-256, синхронизация узлов через Chrony (<1 мс), смена канала через procfs драйвера.
    \item Реализовал \textbf{Bidirectional DSHOT} на \textbf{C} для STM32F7: чтение оборотов моторов без отдельного датчика. \newline \github{RaulIQ/bit-banging-bidir-dshot}
    \item Разработал на \textbf{C} (ESP-IDF, FreeRTOS) двунаправленный UART-мост MAVLink для ESP32-C3/S3 с шифрованием \textbf{AES-128-CTR}: аппаратная и программная реализации, замер времени расшифровки.
  \end{itemize}

  \cvevent{Лаборант Центра БПЛА}{МТУСИ · Март 2024 — Сентябрь 2025}{}{}
  \begin{itemize}
    \item Реализовал протокол \textbf{DSHOT} для управления моторами. \newline \github{RaulIQ/waveform_example}
    \item В команде писал полётную прошивку с нуля: режимы Acro и Stab на \textbf{PID-регуляторах}, обработка данных \textbf{IMU} цифровыми фильтрами, расчёт ориентации на кватернионах.
    \item Разработал систему \textbf{бортовой телеметрии и диагностики} («чёрный ящик»): запись полётных данных на SPI-flash в собственном формате. Добавил неблокирующий async API в форк spi-memory. \newline \github{RaulIQ/spi-memory-async}
  \end{itemize}

  \cvsection{Open Source}
  \cvevent{Контрибьютор Embassy}{Async-фреймворк для embedded Rust, 10 тыс. звёзд}{}{}
  \begin{itemize}
    \item Многоканальный ШИМ через DMA burst (принят): \newline \github{embassy-rs/embassy/pull/4232}
    \item Захват сигнала таймера через DMA, например для DShot (принят): \newline \github{embassy-rs/embassy/pull/4726}
  \end{itemize}

  \cvevent{Контрибьютор open-source библиотек}{}{}{}
  \github{wboayue/fusion-ahrs} \github{sulami/dshot-frame} \newline \github{jettify/uf-ahrs}

  %% Switch to the right column.
  \switchcolumn

  \cvsection{Проекты}

  \cvevent{mik32-rs}{Первая поддержка Rust для RISC-V микроконтроллера MIK32 «Амур»: bring-up по документации, runtime, PAC и HAL}{}{}
  \github{mik32-rs/mik32-rt} \newline \github{mik32-rs/mik32v2-pac} \newline \github{mik32-rs/mik32-hal}

  \medskip

  \cvsection{Образование}

  \textbf{Бакалавриат, МТУСИ} (2022 -- 2027): «Искусственный интеллект и машинное обучение».

  \smallskip
  \textbf{Фокус:} алгоритмы, программная инженерия, ML/AI, DevOps, интеллектуальные системы (CV/SLAM, обработка сигналов).

  \cvsection{Навыки}
  \textbf{Языки:} \cvtag{C++} \cvtag{C} \cvtag{Rust} \newline \cvtag{Python} \\
  \textbf{Платформы:} \cvtag{ESP32} \cvtag{STM32} \cvtag{RPi} \\
  \textbf{Периферия:} \cvtag{DMA} \cvtag{таймеры} \cvtag{прерывания} \\
  \textbf{Среды выполнения:} \cvtag{FreeRTOS} \cvtag{Bare-metal} \newline \cvtag{Embedded Linux} \cvtag{Embassy} \\
  \textbf{Интерфейсы:} \newline \cvtag{I2C/UART/SPI} \cvtag{SBUS/DSHOT} \\
  \textbf{Отладка:} \cvtag{GDB} \cvtag{OpenOCD} \newline \cvtag{JTAG/SWD} \cvtag{осциллограф} \newline \cvtag{лог. анализатор}

  \cvsection{Активности}
  \begin{itemize}
    \item Студкемп «Алиса и умные устройства Яндекса»
    \item \href{https://files.selectel.ru/docs/ru/careerwave_scholarship_finalists.pdf}{Стипендиат Selectel Career Wave}
    \item \href{https://t.me/it_bash_forum/24}{Выступление на форуме IT-Bash}
  \end{itemize}

\end{paracol}

\end{document}
